\chapter{Example: the gate-driven junction}\label{ch:ex-gate}

A microwave voltage on the gate electrode modulates the dot level, $\eps(t)=\eps_0+A_\eps\cos\Omega
t$. Gate driving is the standard way to probe charge transitions in quantum dots and is used
for dispersive gate readout of Andreev and quantum-dot
states~\cite{platero2004photon,petta2004manipulation,wesdorp2023dynamical,sahu2024groundstate}.
In the junction, gate and phase driving are complementary because their coupling
operators obey different selection rules.

\section{Complementary selection rules}

To first order in the amplitude, the resonant splitting of the $S$--$S^*$ transition is
$A_\eps\abs{\matel S{n_d}{S^*}}$ for the gate drive and $A_\phi\abs{\matel S{\partial_\phi
H}{S^*}}$ for the phase drive. \Cref{fig:gate}(a) shows both, computed by the Sambe solver
at the undriven resonance frequency $\Omega=\omega_{SS^*}$, together with the static matrix
elements, as functions of the detuning $\eps_0$. Divided by $A_\eps$, the gate-drive splitting
agrees with $\abs{\matel S{n_d}{S^*}}$ within $\q{gate.gate_deviation}$. Divided by $A_\phi$,
the phase-drive splitting agrees with $\abs{\matel S{\partial_\phi H}{S^*}}$ within $\q{gate.phase_deviation}$, except at
$\eps_0=0$ and at $\eps_0=\pm\q{gate.outlier_position}$. At the latter detunings, levels of the
discretized continuum lie close to the two-photon energy $E_S+2\Omega$ and are coupled to $S^*$
by the phase drive; this shifts the dressed states and changes the splitting per unit amplitude
by $\q{gate.outlier_deviation}$.
The phase-drive matrix element vanishes exactly at the particle--hole-symmetric point
and grows linearly away from it; the gate-drive matrix element is largest there. This is the
eta-parity selection rule of \cref{sec:selection}: $S$ and $S^*$ have opposite eta parity,
which only an eta-odd perturbation such as $n_d-1$ can connect. The large-gap expressions
\eqref{eq:largegap-rabi}, $\propto\sin\vartheta$ and $\propto\cos\vartheta$, reproduce the
trend qualitatively. The small residual phase-drive splitting at $\eps_0=0$ arises from
AC Stark detuning of the resonance while the coupling remains zero.

\begin{physics}
At particle--hole symmetry the singlets $S$ and $S^*$ are, to a good approximation, the
bonding and antibonding combinations of the empty and doubly occupied dot. A phase drive
modulates the strength of the pairing that binds them, shifting their energies while
preserving their distinct eta parities. A gate drive tilts the balance between $\ket0$ and
$\ket2$ and converts $S$ into $S^*$ directly.
\end{physics}

\section{The Floquet spectrum}

\Cref{fig:spectrum}(c,d) is the gate-drive counterpart of the phase-drive fan diagram. Because
the gate drive is linear in $\eps$, it leaves the time-averaged Hamiltonian unchanged
($H_0=H(\eps_0)$). Away from resonances the dressed singlet therefore stays close to $E_S$. The
one-photon $S$--$S^*$ avoided crossing is wide, reflecting the large charge matrix element. The
two-photon crossing is invisible at this amplitude, and the couplings to the continuum
replicas are weak. At fixed $\Omega=\q{spectrum.gate_omega}<\omega_{SS^*}$, panel (d), the
dressed singlet is pushed down by the $S^*$ replica above it. To second order,
$\delta\qe_S\simeq-\tfrac{A_\eps^2}4\abs{\matel S{n_d}{S^*}}^2\bigl[(\omega_{SS^*}-\Omega)^{-1}+(\omega_{SS^*}+\Omega)^{-1}\bigr]$,
the rotating and counter-rotating terms of \cref{eq:ac-stark}.

\section{Two-photon transitions}

At $\Omega\approx\omega_{SS^*}/2$ the pair transition is driven by two photons.
\Cref{fig:gate}(b) compares the one- and two-photon splittings at their respective
resonances. The one-photon splitting grows linearly with $A_\eps$, the two-photon splitting
quadratically, as a second-order process must. Two mechanisms contribute to the two-photon
coupling: virtual transitions through other states, and the difference of the average
charges of $S$ and $S^*$, which modulates their splitting and acts like the Bessel-function
mechanism of \cref{ch:ex-tls}. Both resonances move with the drive amplitude, by AC Stark and
Bloch--Siegert shifts; each resonance is located at the minimum of the quasienergy splitting
as a function of $\Omega$. At the largest amplitudes the one-photon splitting falls below the
linear dependence as counter-rotating terms and coupling to other states become important.

\begin{figure}[htb]
\centering
\manualfigure{gate}
\caption{Gate-driven reference junction. (a) Resonant $S$--$S^*$ splitting per unit amplitude
for a gate drive ($A_\eps=\q{gate.gate_amplitude}$, circles) and a phase drive
($A_\phi=\q{gate.phase_amplitude}$, squares, scaled by $\q{gate.phase_scale}$), versus the
detuning $\eps_0$, compared with the static matrix elements (lines). (b) One- and two-photon
splittings at their resonances versus the gate amplitude; dotted lines are $\propto A_\eps$
and $\propto A_\eps^2$.}
\label{fig:gate}
\end{figure}

\begin{incode}
\begin{lstlisting}[frame=none,backgroundcolor={},aboveskip=0pt,belowskip=0pt]
gate = PeriodicHamiltonian.from_parameters(reference_model, model,
                                           ParameterDrive.cosine("detuning", .01), omega=.716)
states = solve_floquet(gate, 2, Sector(0, 0), dict(floquet_cutoff=2, states=2))
\end{lstlisting}
Since a linear gate drive has only $H_{\pm1}$, a few samples reproduce $H(t)$ exactly; the
two-photon calculation needs $M\ge2$ to contain the intermediate replica.
\par\smallskip\usermanual{docs/floquet.md}.
\end{incode}

\tutorial{07_gate_drive}
